\begin{figure}[ht]
\centering
\begin{tikzpicture}[scale=1.0,
  bar/.style = {draw=engblack, thick},
]
% Three stacked-bar cartoons: work per level down the recursion tree.
% Case 1 (leaf-heavy): bars widen going down.
% Case 2 (balanced): bars equal width.
% Case 3 (root-heavy): bars narrow going down.
\def\lvls{0,1,2,3}

% Case 1: leaves dominate (widths grow).
\begin{scope}[shift={(0,0)}]
  \node[font=\small, anchor=south] at (1.1, 3.4) {Case 1: leaves};
  \foreach \k/\w in {0/0.5, 1/0.9, 2/1.5, 3/2.2} {
    \fill[engred!30] (1.1-\w/2, 2.8-\k*0.7) rectangle (1.1+\w/2, 2.8-\k*0.7+0.45);
    \draw[bar]       (1.1-\w/2, 2.8-\k*0.7) rectangle (1.1+\w/2, 2.8-\k*0.7+0.45);
  }
  \node[font=\scriptsize, anchor=north, text width=2.6cm, align=center]
    at (1.1, 0.35) {$f(n)=O(n^{\log_b a-\epsilon})$\\ $T=\Theta(n^{\log_b a})$};
\end{scope}

% Case 2: balanced (widths equal).
\begin{scope}[shift={(4.0,0)}]
  \node[font=\small, anchor=south] at (1.1, 3.4) {Case 2: balanced};
  \foreach \k in {0,1,2,3} {
    \fill[englime!40] (1.1-0.75, 2.8-\k*0.7) rectangle (1.1+0.75, 2.8-\k*0.7+0.45);
    \draw[bar]        (1.1-0.75, 2.8-\k*0.7) rectangle (1.1+0.75, 2.8-\k*0.7+0.45);
  }
  \node[font=\scriptsize, anchor=north, text width=2.6cm, align=center]
    at (1.1, 0.35) {$f(n)=\Theta(n^{\log_b a})$\\ $T=\Theta(n^{\log_b a}\log n)$};
\end{scope}

% Case 3: root dominates (widths shrink).
\begin{scope}[shift={(8.0,0)}]
  \node[font=\small, anchor=south] at (1.1, 3.4) {Case 3: root};
  \foreach \k/\w in {0/2.2, 1/1.5, 2/0.9, 3/0.5} {
    \fill[engblue!25] (1.1-\w/2, 2.8-\k*0.7) rectangle (1.1+\w/2, 2.8-\k*0.7+0.45);
    \draw[bar]        (1.1-\w/2, 2.8-\k*0.7) rectangle (1.1+\w/2, 2.8-\k*0.7+0.45);
  }
  \node[font=\scriptsize, anchor=north, text width=2.6cm, align=center]
    at (1.1, 0.35) {$f(n)=\Omega(n^{\log_b a+\epsilon})$\\ $T=\Theta(f(n))$};
\end{scope}
\end{tikzpicture}
\caption{The three regimes of the master theorem, read as the
distribution of work across the levels of the recursion tree. Each
bar is the total work at one level, the root at top and the leaves
at bottom. In case 1 the work grows toward the leaves, so the leaf
term $n^{\log_b a}$ dominates; in case 2 every level carries the
same work and the total picks up a factor $\log n$ from the level
count; in case 3 the root's work dominates and swamps the geometric
tail below it. Reading which way the bars widen is the fastest way to
name the case before touching the algebra.}
\label{fig:vol02:ch17:master-regimes}
\end{figure}
